\chapter{Euclid幾何学}
    \section{諸公式}
        \subsection{ヘロンの公式}
            \begin{shadebox}
                辺の長さが$a,b,c$なる三角形の面積$S$は
                \[S = \sqrt{s(s-a)(s-b)(s-c)}\]
                但し
                \[s = \frac{a+b+c}{2}\]
            \end{shadebox}
            \begin{proof}
                余弦定理を用いて
                \[\cos{C} = \frac{a^2+b^2-c^2}{2ab}\]
                \[\therefore \; \sin{C} = \sqrt{1-\cos^2{C}} = \frac{\sqrt{4a^2b^2-{\parens*{a^2+b^2-c^2}}^2}}{2ab}\]
                \begin{align*}
                    S &=& \frac{1}{2}ab\sin{C}\\
                    &=& \frac{1}{2}\sqrt{4a^2b^2-{\parens*{a^2+b^2-c^2}}^2}
                \end{align*}
                ここで
                \begin{align*}
                    &\;& 4a^2b^2-{\parens*{a^2+b^2-c^2}}^2\\
                    &=& -\parens*{a^4+b^4+c^4-2a^2b^2-2b^2c^2-2c^2a^2}\\
                    &=& -\bracks*{a^4 - 2(b^2+c^2)a^2 + b^4+c^4 -2b^2c^2}\\
                    &=& -\bracks*{a^4 - 2(b^2+c^2)a^2 + (b^2+c^2)^2 -4b^2c^2}\\
                    &=& -\braces*{\bracks*{a^2-(b^2+c^2)^2}^2-(2bc)^2}\\
                    &=& -\bracks*{a^2-(b^2+c^2)+2bc}\bracks*{a^2-(b^2+c^2)-2bc}\\
                    &=& -\bracks*{a^2-(b-c)^2}\bracks*{a^2-(b+c)^2}\\
                    &=& (a+b+c)(-a+b+c)(a-b+c)(a+b-c)
                \end{align*}
                よって
                \begin{align*}
                    S &=& \sqrt{\frac{a+b+c}{2}\times\frac{-a+b+c}{2}\times\frac{a-b+c}{2}\times\frac{a+b-c}{2}}\\
                    &=& \sqrt{s(s-a)(s-b)(s-c)}
                \end{align*}
                但し
                \[s = \frac{a+b+c}{2}\]
            \end{proof}
    \section{諸定理}
        \subsection{\texorpdfstring{$\lim_{\norm{\vx}\to\infty}\norm{\vx-\bm{a}} - \norm{\vx} = -\frac{\vx}{\norm{\vx}}\cdot\bm{a}$}{近接する2点と遠方の点との距離の差の極限}}
            \begin{shadebox}
                $\bm{a},\vx\in\realNumbers^3$とするとき次式が成り立つ。
                \[ \lim_{\norm{\vx}\to\infty}\norm{\vx-\bm{a}} - \norm{\vx} = -\frac{\vx}{\norm{\vx}}\cdot\bm{a} \]
            \end{shadebox}
            \begin{proof}
                \quad\par
                $\bm{a}=\bm{0}$のときは明らかに成り立つ。
                以下では$\bm{a}\neq 0$とする。
                % ↓ 没 @2025-01-12
                % また，$\norm{\vx}$を十分大きくとり，$\norm{\vx}>\norm{\bm{a}}/2$とする。
                % $f(\vx) := \norm{\vx-\bm{a}} - \norm{\vx}$とすると次式が成り立つ。
                % \[ \parens*{f(\vx) + \norm{\vx}}^2 = \norm{\vx-\bm{a}}^2 \quad \therefore f(\vx)\parens*{f(\vx) + 2\norm{\vx}} = \norm{\bm{a}}^2 - 2\vx\cdot\bm{a} \]
                % $f(\vx)$が最小となるのは，原点と$\vx$を結ぶ線分の上に$\bm{a}$があるときで，その値は$-\norm{\bm{a}}$である。
                % $\norm{\vx}>\norm{\bm{a}}/2$と仮定しているから次式を得る。
                % \[ f(\vx) = \frac{\norm{\bm{a}}^2 - 2\vx\cdot\bm{a}}{f(\vx) + 2\norm{\vx}} < \frac{\norm{\bm{a}}^2 - 2\vx\cdot\bm{a}}{2\norm{\vx} - \norm{\bm{a}}} \to -\frac{\vx}{\norm{\vx}}\cdot\bm{a} \quad \text{as} \quad \norm{\vx}\to\infty\]
                % ↓ new @2025-01-12
                \begin{align*}
                    % \norm{\vx-\bm{a}} - \norm{\vx} &= \frac{(\norm{\vx-\bm{a}} - \norm{\vx})(\norm{\vx-\bm{a}} + \norm{\vx})}{\norm{\vx-\bm{a}} + \norm{\vx}} = \frac{\norm{\vx-\bm{a}}^2 - \norm{\vx}^2}{\norm{\vx-\bm{a}} + \norm{\vx}} \\
                    % &= \frac{\norm{\bm{a}}^2 - 2\vx\cdot\bm{a}}{\norm{\vx-\bm{a}} + \norm{\vx}}
                    \norm{\vx-\bm{a}} - \norm{\vx} &= \frac{(\norm{\vx-\bm{a}} - \norm{\vx})(\norm{\vx-\bm{a}} + \norm{\vx})}{\norm{\vx-\bm{a}} + \norm{\vx}} = \frac{\norm{\vx-\bm{a}}^2 - \norm{\vx}^2}{\norm{\vx-\bm{a}} + \norm{\vx}} \\
                    &= \frac{\norm{\bm{a}}^2 - 2\vx\cdot\bm{a}}{\norm{\vx-\bm{a}} + \norm{\vx}} = \frac{\norm{\bm{a}}^2/\norm{\vx} - 2\vx\cdot\bm{a}/\norm{\vx}}{1+\frac{\norm{\vx-\bm{a}}}{\norm{\vx}}} \to -\frac{\vx}{\norm{\vx}}\cdot\bm{a} \as \norm{\vx}\to\infty
                \end{align*}
            \end{proof}